\chapter{テキストと LLM}
\label{ch:text}

LLM-BLM の研究に最も近い分野である。ニュース、決算発表、決算説明会、提出書類の文章から、株価の予測に使える情報を取り出す研究は、辞書から機械学習、そして LLM へと進んできた。

\section{辞書の時代}

\begin{paperbox}[Tetlock（2007）\cite{tetlock2007}]
Wall Street Journal の株式市場のコラムの文章の悲観の度合い（一般の心理の辞書で数えた否定的な語の割合）が高い日の翌日は、株価が下がり、その後数日で戻ることを示した。悲観の高さは、売買高の多さも予測した。メディアの語調が、投資家の心理を通じて、一時的に価格を動かすという解釈である。
\end{paperbox}

Tetlock・Saar-Tsechansky・Macskassy（2008）\cite{tsm2008} は、個別の会社についてのニュースの否定的な語の割合が、その会社の将来の利益とリターンの低さを予測し、市場はその情報にわずかに遅れて反応することを示した。

\begin{paperbox}[Loughran・McDonald（2011）\cite{lm2011}]
一般の心理の辞書（Harvard の辞書）で「否定的」とされる語の約 4 分の 3 は、10-K（米国の年次報告書）の文脈では否定的な意味を持たない（「負債」「税」「費用」など）ことを示し、金融の文章のための辞書を作った。この辞書による否定的な語の割合は、提出の前後の株価の反応などと関係する。
\end{paperbox}

LLM-BLM の \code{dictionary} の点数は、この Loughran--McDonald の辞書による語調である。

\subsection{文章の変化}

\begin{paperbox}[Cohen・Malloy・Nguyen（2020）「怠惰な価格」\cite{cmn2020}]
10-K と 10-Q の文章が、前年の同じ書類から変わった度合いを測り、文章を大きく変えた会社（変更者）の株が、その後大きく下がることを示した。変更者を売り、変えなかった会社を買うポートフォリオは、月に最大 188 ベーシスポイント（年 22\% 超）のアルファを上げた。投資家は長い書類の変化を読まず、情報は数か月かけて価格に入る。
\end{paperbox}

この研究は、「文章の情報は、数値の情報よりゆっくり価格に入る」ことの強い証拠の一つで、PEAD.txt（第 \ref{ch:event} 章）と同じ方向の結果である。

\section{機械学習の時代}

\begin{itemize}
\item \textbf{Ke・Kelly・Xiu（2019）}\cite{kkx2019}：ニュースの文章から、リターンの符号と一緒に現れやすい語を教師ありで学び、語調の点数を作る方法（SESTM）を提案した。この点数は、既存の商用の語調の指標より強くリターンを予測し、その情報は、小さく変動の大きい株ほどゆっくり価格に入った。
\item \textbf{Frankel・Jennings・Lee（2022）}\cite{fjl2022}：決算説明会や 10-K の文章で、機械学習による語調の測り方が、辞書による測り方より、開示に対する株価の反応をよく説明することを示した。
\item \textbf{FinBERT}：言語モデル BERT を金融の文章で追加学習したもの。Araci（2019）\cite{araci2019} は金融のニュースの語調の分類で、Huang・Wang・Yang（2023）\cite{hwy2023} はアナリストの報告書などで、辞書や従来の機械学習より高い精度を示した。LLM-BLM の \code{finbert} と \code{finbert\_ft} はこの系統である。
\end{itemize}

\section{LLM の時代}

\subsection{LLM はリターンを予測するか}

\begin{paperbox}[Lopez-Lira・Tang（2023）\cite{lopezliratang2023}]
ChatGPT に、ニュースの見出しがその会社の株価に良いか悪いかを答えさせた点数が、その後のリターンを予測することを示した。LLM の言語の理解の能力が、金融の予測の能力として現れたことを示した最初期の研究の一つである。改訂を重ねた arXiv の最新の版の要旨によれば、
\begin{itemize}
\item GPT-4 は、学習の知識の期限より後の見出しで、取引できない最初の反応（見出しが出た時点の反応）の方向を、ポートフォリオの日の約 90\% で当てた。
\item GPT-4 の点数は、その後のドリフトも有意に予測し、特に小さな株と悪いニュースで強かった。
\item 予測の力はモデルが大きいほど高い。
\item LLM の利用が広がるにつれて、戦略の収益は下がった（価格の効率が上がったことと整合的）。
\end{itemize}
「約 90\%」は取引できない最初の反応についての値で、取引できる収益の値ではないことに注意する（Deep Research の答え 2 は、この区別を落としていた）。
\end{paperbox}

\begin{paperbox}[Chen・Kelly・Xiu（2022）「期待リターンと LLM」\cite{ckx2022}]
LLM の点数（語調）ではなく、LLM が文章から作るベクトル（埋め込み）全体を説明変数にして、ニュースからリターンを予測した。世界の多くの市場で、埋め込みによる予測は、語調の点数や従来の言語の方法より強く、大きなモデルほど予測が良かった。
\end{paperbox}

\begin{paperbox}[Kim・Muhn・Nikolaev（2024）\cite{kmn2024}]
GPT-4 に、会社名を伏せた財務諸表（貸借対照表と損益計算書）だけを与えて、来期の利益が増えるか減るかを予測させた。GPT-4 の予測の正しさは、アナリストの予想を上回り、専門の機械学習のモデルに匹敵した。その予測による売買の戦略は、高いシャープレシオとアルファを示した。
\end{paperbox}

\subsection{先読みの問題}

LLM の研究の最大の問題は、\textbf{先読み}（\chapref{ch:backtest}）である。LLM は、学習したデータの期間の出来事を知っている。2015 年の決算発表文を、2023 年までの文章で学習した LLM に読ませると、LLM はその会社のその後の株価や業績を「知っている」可能性がある。

\begin{itemize}
\item \textbf{Glasserman・Lin（2023）}\cite{glasserman2023}：GPT の語調の点数による売買の成績を、元の見出しと、会社名などを伏せた見出しで比べた。LLM の学習の期間の中では、意外にも伏せた見出しのほうが成績が良く、特に大きな会社で強かった。会社についての一般の知識が語調の判断を乱す効果（気をそらす効果）のほうが、先読みより大きかったと解釈した。学習の期間の外では、先読みは問題にならないが、気をそらす効果は残りうる。
\item \textbf{ChronoGPT}（He・Lv・Manela・Wu 2025）\cite{chronogpt2025}：各時点までに手に入った文章だけで学習した言語モデルの系列（ChronoBERT・ChronoGPT）を作り、先読みのない検証を可能にした。ニュースから翌日のリターンを予測する応用では、その場その場の出力によるシャープレシオが、ずっと大きな Llama のモデルと同程度で、先読みの偏りは小さいと報告した。LLM-BLM の主な検定は、2015 年までの文章で学習した版を使った。
\end{itemize}

\subsection{金融に特化した LLM}

\begin{itemize}
\item \textbf{BloombergGPT}（Wu ほか 2023）\cite{bloomberggpt2023}：Bloomberg が、自社の金融のデータ（3,630 億トークン）と一般のデータ（3,450 億トークン）で、500 億のパラメータの LLM を一から学習した。金融の課題では、既存のモデルより大きく良く、一般の課題でも劣らなかった、と報告した。
\item \textbf{FinGPT}（Yang・Liu・Wang 2023）\cite{fingpt2023}：公開の LLM を、金融のデータで安く追加学習する（LoRA などの軽い方法）枠組みを公開した。一から学習するより、はるかに安い。
\end{itemize}

\section{効果は残っているか}

文章の情報の予測力については、次のことが分かっている。

\begin{enumerate}
\item \textbf{発表の直後の反応}は、文章の点数で強く説明できる（LLM-BLM の \SP{} の実験もこれを確認した）。
\item \textbf{翌日以降のドリフト}は、大型株では小さいか、消えている（第 \ref{ch:event} 章の Martineau 2022）。残っているのは、小さな株、悪い知らせ、長く複雑な文章（10-K の変化、決算説明会の発言）である。
\item \textbf{LLM の普及}は、文章の情報の処理を安く速くし、こうした遅れを小さくしていくと予想される。過去のデータで見つかった LLM の戦略の成績が、これからも続くとは限らない。
\item \textbf{先読み}を除くと、学習の期間の中での LLM の成績は、小さくなりうる。
\end{enumerate}

\begin{projbox}
LLM-BLM の結果（2 日の反応は強く説明、ドリフトは有意でない）は、この分野の研究と整合的である。研究として新しい情報を加えられる余地は、(1) 先読みのない LLM（ChronoGPT）での検証、(2) 注目の低い小型株でのドリフト（実験 S1）、(3) 数値でなく文章でしか分からない情報（ガイダンス、リスクの記述、文章の変化）、(4) 発表の当日の反応の大きさの予測（取引ではなく、リスク管理やイベントの前の量の調整に使う）、などにある。
\end{projbox}
